La clase anterior estableció la dualidad en categorías monoidales simétricas y la relación entre tener un dual y una adjunción. Esta clase retoma esa discusión desde dos frentes. Por un lado se extraen las consecuencias formales de la dualidad —la unidad es autodual, el dual de un producto tensorial es el producto tensorial de los duales, el dual es único salvo isomorfismo— y se aísla el teorema de coherencia que permite abreviar todos los cálculos diagramáticos. Por otro lado se introduce el \emph{hom interno}, que es la pieza que en una categoría monoidal simétrica general ocupa el lugar del dual, y se estudia la noción de categoría monoidal simétrica \emph{cerrada}.

La segunda mitad de la clase vuelve al contexto estable. La categoría de Spanier--Whitehead con el producto smash no se sabe todavía que sea cerrada, de modo que no hay un candidato evidente para el dual de un objeto arbitrario; pero las esferas sí tienen dual, y se demuestra que basta con comprobar la dualidad sobre las esferas para obtenerla en toda la categoría. El caso general se deja para la clase siguiente.

\begin{nota}
   La clase se abre reenunciando el criterio de la sesión anterior, que se usará una y otra vez: \emph{$X^\vee$ es un dual de $X$ si y solo si $-\otimes X$ es adjunto izquierdo a $-\otimes X^\vee$}. Es el enunciado de la \myref{prop:adjuncion} leído como una equivalencia entre la existencia del dual y la existencia de la adjunción, y es la forma en que se aplicará en todo el capítulo.
\end{nota}

%%%%%%%%%%%%%%%%%%%%%%%%%%%%%%%%%%%%%%%%%%%%%%%%%%%%%%%%%%%%%%%%%%%%%%%%%%%%%%%
\section{Consecuencias formales de la dualidad}
%%%%%%%%%%%%%%%%%%%%%%%%%%%%%%%%%%%%%%%%%%%%%%%%%%%%%%%%%%%%%%%%%%%%%%%%%%%%%%%

%%%%%%%%%%%%%%%%%%%%%%%%%%%%%%%%%%%%%%%%%%%%%%%%%%%%%%%%%%%%%%%%%%%%%%%%%%%%%%%
\subsection{Corrección del diagrama de la clase anterior}
%%%%%%%%%%%%%%%%%%%%%%%%%%%%%%%%%%%%%%%%%%%%%%%%%%%%%%%%%%%%%%%%%%%%%%%%%%%%%%%

La clase comienza con una corrección. Al final de la clase anterior se escribió un diagrama que debía probar que las transformaciones $\Phi$ y $\Psi$ de la \myref{prop:adjuncion} son inversas, y el diagrama presentado era ya una \emph{simplificación}: faltaba el paso que reduce el cálculo a ese diagrama. El paso omitido es un cuadrado de naturalidad, y su ausencia hacía que el término de una de las esquinas apareciera desplazado.

La versión completa, siguiendo los pasos en el orden en que se construyen, es la siguiente. Partimos de $f\colon A\to B\otimes X^\vee$ y aplicamos $\Psi$ y después $\Phi$: la primera parte del cálculo es
\[
\begin{aligned}
A
&\xrightarrow{\ \rho_A^{-1}\ } A\otimes1
\xrightarrow{\ 1_A\otimes\mathrm{coev}\ } A\otimes(X\otimes X^\vee)
\xrightarrow{\ \alpha^{-1}_{A,X,X^\vee}\ } (A\otimes X)\otimes X^\vee ,
\end{aligned}
\]
y a partir de ahí el cálculo se abre en dos caminos que deben coincidir. El primero pasa por el espacio de la izquierda,
\[
\begin{aligned}
(A\otimes X)\otimes X^\vee
&\xrightarrow{\ (f\otimes1_X)\otimes1_{X^\vee}\ } ((B\otimes X^\vee)\otimes X)\otimes X^\vee
\xrightarrow{\ \alpha_{B,X^\vee,X}\otimes1_{X^\vee}\ } (B\otimes(X^\vee\otimes X))\otimes X^\vee\\
&\xrightarrow{\ (1_B\otimes\mathrm{ev})\otimes1_{X^\vee}\ } (B\otimes1)\otimes X^\vee
\xrightarrow{\ \rho_B\otimes1_{X^\vee}\ } B\otimes X^\vee ,
\end{aligned}
\]
y el segundo reasocia primero y elimina las unidades al final,
\[
\begin{aligned}
A\otimes(X\otimes X^\vee)
&\xrightarrow{\ f\otimes1_{X\otimes X^\vee}\ } (B\otimes X^\vee)\otimes(X\otimes X^\vee)
\xrightarrow{\ \alpha^{-1}_{B,X^\vee,X\otimes X^\vee}\ } B\otimes(X^\vee\otimes(X\otimes X^\vee))\\
&\xrightarrow{\ 1_B\otimes \alpha^{-1}_{X^\vee,X,X^\vee}\ } B\otimes((X^\vee\otimes X)\otimes X^\vee)
\xrightarrow{\ 1_B\otimes(\mathrm{ev}\otimes1_{X^\vee})\ } B\otimes(1\otimes X^\vee)\\
&\xrightarrow{\ 1_B\otimes\lambda_{X^\vee}\ } B\otimes X^\vee .
\end{aligned}
\]
Los dos caminos no empiezan en el mismo sitio: el primero en $(A\otimes X)\otimes X^\vee$ y el segundo en $A\otimes(X\otimes X^\vee)$. Lo que los conecta es precisamente el cuadrado de naturalidad del asociador
\[
\begin{tikzcd}[column sep=large, row sep=large]
(A\otimes X)\otimes X^\vee \arrow[r,"\alpha_{A,X,X^\vee}"] \arrow[d,"(f\otimes1_X)\otimes1_{X^\vee}"'] & A\otimes(X\otimes X^\vee) \arrow[d,"f\otimes1_{X\otimes X^\vee}"]\\
((B\otimes X^\vee)\otimes X)\otimes X^\vee \arrow[r,"\alpha_{B\otimes X^\vee,X,X^\vee}"'] & (B\otimes X^\vee)\otimes(X\otimes X^\vee)
\end{tikzcd}
\]
que es el paso que se había omitido. En la pizarra del día anterior faltaba esta esquina, y por eso el término de la derecha aparecía escrito en el lugar del de la izquierda.

\begin{nota}
   La demostración que aparece en la \myref{prop:adjuncion} es exactamente la descomposición de este diagrama en sus piezas: el cuadrado de naturalidad anterior, el pentágono para $(B,X^\vee,X,X^\vee)$, la eliminación de las unidades y, para cerrar, la identidad de zigzag junto con el \myref{lema:coherencia}. Por eso la corrección no obliga a cambiar nada de lo ya escrito; se deja constancia de ella porque el diagrama completo, tal como se escribió en la pizarra, es el que muestra de una sola vez que la composición vale $f$.

   Conviene precisar qué había exactamente en la pizarra: allí el diagrama consta del cuadrado de naturalidad del asociador, del tramo horizontal $A\otimes(X\otimes X^\vee)\to(B\otimes X^\vee)\otimes(X\otimes X^\vee)$ y de \emph{una sola} columna de descenso, la que baja desde $((B\otimes X^\vee)\otimes X)\otimes X^\vee$. Del vértice $(B\otimes X^\vee)\otimes(X\otimes X^\vee)$ no sale ninguna flecha, y la convergencia de los dos caminos se lee sobre la columna compartida. El «segundo camino» escrito aquí con sus cinco pasos hace explícita esa lectura.
\end{nota}

%%%%%%%%%%%%%%%%%%%%%%%%%%%%%%%%%%%%%%%%%%%%%%%%%%%%%%%%%%%%%%%%%%%%%%%%%%%%%%%
\subsection{Consecuencias inmediatas: la unidad y el producto tensorial}
%%%%%%%%%%%%%%%%%%%%%%%%%%%%%%%%%%%%%%%%%%%%%%%%%%%%%%%%%%%%%%%%%%%%%%%%%%%%%%%

De la caracterización de la dualidad por adjunciones se obtienen dos consecuencias elementales que se usarán constantemente.

\begin{coro}\label{coro:unidad-autodual}
La unidad monoidal $1$ es un dual de sí misma.
\end{coro}
\begin{proof}
Por la \myref{prop:adjuncion}, basta exhibir un isomorfismo natural entre los funtores $-\otimes1$ y $-\otimes1$. Pero
\[
\hom_{\mathcal C}(A\otimes1,\ B)\cong\hom_{\mathcal C}(A,\ B\otimes1)
\]
naturalmente en $A$ y en $B$: ambos miembros son naturalmente isomorfos a $\hom_{\mathcal C}(A,B)$, porque $A\otimes1\cong A$ y $B\otimes1\cong B$ mediante las unidades $\rho_A$ y $\rho_B$, que son naturales. Los dos funtores son por tanto el mismo y la adjunción es la identidad.
\end{proof}

\begin{coro}\label{coro:dual-tensor}
Si $X^\vee$ e $Y^\vee$ son duales de $X$ e $Y$ respectivamente, entonces $Y^\vee\otimes X^\vee$ es un dual de $X\otimes Y$.
\end{coro}
\begin{proof}
Se aplica la caracterización de la dualidad dos veces, omitiendo los paréntesis por comodidad ya que los asociadores son naturales:
\[
\hom_{\mathcal C}(A\otimes X\otimes Y,\ B)
\cong\hom_{\mathcal C}(A\otimes X,\ B\otimes Y^\vee)
\cong\hom_{\mathcal C}(A,\ B\otimes Y^\vee\otimes X^\vee),
\]
donde el primer isomorfismo usa que $Y^\vee$ es un dual de $Y$ y el segundo que $X^\vee$ lo es de $X$. La composición es un isomorfismo natural en $A$ y en $B$, de modo que $-\otimes(X\otimes Y)$ es adjunto izquierdo de $-\otimes(Y^\vee\otimes X^\vee)$, es decir, $Y^\vee\otimes X^\vee$ es un dual de $X\otimes Y$.
\end{proof}

Obsérvese el orden: el dual de $X\otimes Y$ es $Y^\vee\otimes X^\vee$ y no $X^\vee\otimes Y^\vee$. El orden es el esperado para un funtor contravariante, y en una categoría simétrica ambos coinciden salvo la trenza, que no es la identidad pero sí un isomorfismo.

\begin{defi}\label{defi:dualizable}
Un objeto $X$ de una categoría monoidal simétrica es \emph{dualizable} si admite un dual.
\end{defi}

%%%%%%%%%%%%%%%%%%%%%%%%%%%%%%%%%%%%%%%%%%%%%%%%%%%%%%%%%%%%%%%%%%%%%%%%%%%%%%%
\subsection{El teorema de coherencia de Mac Lane y Kelly}
%%%%%%%%%%%%%%%%%%%%%%%%%%%%%%%%%%%%%%%%%%%%%%%%%%%%%%%%%%%%%%%%%%%%%%%%%%%%%%%

Los cálculos de la clase anterior y de esta son diagramas formados por asociadores, unidades y trenzas. El resultado siguiente dice que no hace falta verificarlos uno por uno.

\begin{teo}[Mac Lane--Kelly]\label{teo:mackane-kelly}
Sea $(\mathcal C,\otimes,1,a,\lambda,\rho,B)$ una categoría con un producto monoidal, una unidad, un asociador $\alpha$, unidades $\lambda$ y $\rho$ y una trenza $B$. Entonces $(\mathcal C,\otimes,1,a,\lambda,\rho,B)$ es una categoría monoidal simétrica si y solo si todo diagrama formado por tensores finitos de $a$, $\lambda$, $\rho$ y $B$ es conmutativo.
\end{teo}

\begin{nota}
   La dirección no trivial es la que se usa en la práctica: si $(\mathcal C,\otimes,1,a,\lambda,\rho,B)$ es simétrica monoidal, entonces cualquier diagrama cuyas aristas se obtengan tensando y componiendo $a$, $\lambda$, $\rho$ y $B$ —y las identidades y los morfismos dados— conmuta. Los diagramas que aparecen en la \myref{defi:monoidal} son algunos de estos, y en este sentido los axiomas \emph{generan} todos los demás. La demostración se debe a Mac Lane para el caso de la asociatividad y a Mac Lane y Kelly para el caso con la trenza (véanse \cite{maclane1998} y \cite{kelly1964}).

   En lo que sigue se invocará este teorema para justificar de una vez los reordenamientos que solo involucran $a$, $\lambda$, $\rho$ y $B$; cuando el diagrama incluya además morfismos arbitrarios, habrá que usar su naturalidad.
\end{nota}

%%%%%%%%%%%%%%%%%%%%%%%%%%%%%%%%%%%%%%%%%%%%%%%%%%%%%%%%%%%%%%%%%%%%%%%%%%%%%%%
\subsection{La dualidad es simétrica}
%%%%%%%%%%%%%%%%%%%%%%%%%%%%%%%%%%%%%%%%%%%%%%%%%%%%%%%%%%%%%%%%%%%%%%%%%%%%%%%

La primera consecuencia formal de la definición es que la relación «ser dual de» es simétrica.

\begin{prop}\label{prop:simetria-dual}
Si $X^\vee$ es un dual de $X$, entonces $X$ es un dual de $X^\vee$.
\end{prop}

\begin{proof}
Sean
\[
\mathrm{ev}\colon X^\vee\otimes X\longrightarrow1,
\qquad
\mathrm{coev}\colon1\longrightarrow X\otimes X^\vee
\]
los morfismos de la dualidad. Se definen
\[
\mathrm{ev}'\coloneqq\mathrm{ev}\circ B_{X,X^\vee}\colon X\otimes X^\vee\longrightarrow X^\vee\otimes X\longrightarrow1,
\qquad
\mathrm{coev}'\coloneqq B_{X,X^\vee}\circ\mathrm{coev}\colon1\longrightarrow X\otimes X^\vee\longrightarrow X^\vee\otimes X,
\]
es decir, se permutan los factores con la trenza y después se evalúa, y dualmente se coevalúa y después se permutan:
\[
\begin{tikzcd}[column sep=large, row sep=large]
X\otimes X^\vee \arrow[r,"B_{X,X^\vee}"] \arrow[dr,"\mathrm{ev}'"'] & X^\vee\otimes X \arrow[d,"\mathrm{ev}"]\\
& 1
\end{tikzcd}
\qquad\qquad
\begin{tikzcd}[column sep=large, row sep=large]
1 \arrow[r,"\mathrm{coev}"] \arrow[dr,"\mathrm{coev}'"'] & X\otimes X^\vee \arrow[d,"B_{X,X^\vee}"]\\
& X^\vee\otimes X
\end{tikzcd}
\]
La dirección en que se aplica la trenza es indiferente. La dualidad en una categoría monoidal simétrica es el análogo de la noción de anillo conmutativo, donde el orden de los factores de un producto no importa, y aquí tampoco: las dos trenzas $B_{X,X^\vee}$ y $B_{X^\vee,X}$ son inversas entre sí, de modo que permutar los factores en un sentido o en el otro da el mismo resultado. Por eso la elección anterior es la natural, y la que hace encajar los dominios y codominios de $\mathrm{ev}'$ y $\mathrm{coev}'$.

Para que $X$ sea un dual de $X^\vee$ hay que comprobar las dos identidades de zigzag con $\mathrm{ev}'$ y $\mathrm{coev}'$; por ejemplo, la que afirma que la composición
\[
X^\vee\xrightarrow{\ \lambda_{X^\vee}^{-1}\ }1\otimes X^\vee\xrightarrow{\ \mathrm{coev}'\otimes1_{X^\vee}\ }(X^\vee\otimes X)\otimes X^\vee\xrightarrow{\ \alpha_{X^\vee,X,X^\vee}\ }X^\vee\otimes(X\otimes X^\vee)\xrightarrow{\ 1_{X^\vee}\otimes\mathrm{ev}'\ }X^\vee\otimes1\xrightarrow{\ \rho_{X^\vee}\ }X^\vee
\]
es $1_{X^\vee}$. Obsérvese que $\mathrm{coev}'$ aterriza en $X^\vee\otimes X$ y $\mathrm{ev}'$ sale de $X\otimes X^\vee$, que es lo que hace que las componentes encajen. La verificación se hace sustituyendo las definiciones y reordenando. Como todos los morfismos que intervienen en los reordenamientos son $a$, $\lambda$, $\rho$ y $B$, el \myref{teo:mackane-kelly} permite trabajar en el modelo estricto, donde $a$, $\lambda$ y $\rho$ son identidades; allí la composición anterior se reduce a
\[
(1_{X^\vee}\otimes\mathrm{ev}')\circ(\mathrm{coev}'\otimes1_{X^\vee}),
\]
y sustituyendo $\mathrm{coev}'=B_{X,X^\vee}\circ\mathrm{coev}$ y $\mathrm{ev}'=\mathrm{ev}\circ B_{X,X^\vee}$ esta expresión es la identidad de zigzag original de $X^\vee$ ---que vale $1_{X^\vee}$ por hipótesis--- con la trenza insertada dos veces:
\[
(1_{X^\vee}\otimes\mathrm{ev}')\circ(\mathrm{coev}'\otimes1_{X^\vee})=(1_{X^\vee}\otimes\mathrm{ev})\circ(1_{X^\vee}\otimes B_{X,X^\vee})\circ(B_{X,X^\vee}\otimes1_{X^\vee})\circ(\mathrm{coev}\otimes1_{X^\vee}).
\]
Las dos inserciones de $B_{X,X^\vee}$ vienen de las definiciones de $\mathrm{coev}'$ y $\mathrm{ev}'$: la primera del producto $\mathrm{coev}'\otimes1_{X^\vee}$ y la segunda del producto $1_{X^\vee}\otimes\mathrm{ev}'$. El axioma del hexágono, que es el que gobierna cómo interactúan las trenzas con la composición, permite reordenarlas y reducir la expresión a la identidad de zigzag de la hipótesis. En el caso de los espacios vectoriales de dimensión finita el cálculo es explícito: tomando $\mathrm{coev}(1)=\sum_i v_i\otimes\varphi^i$ se obtiene
\[
\sum_i\varphi^i\otimes v_i\otimes\varphi
\ \longmapsto\ \sum_i\varphi^i\otimes\varphi\otimes v_i
\ \longmapsto\ \sum_i\varphi^i\,\varphi(v_i)=\varphi,
\]
es decir, $1_{X^\vee}$. La otra identidad de zigzag se prueba de manera parecida.

\begin{nota}
   Este paso es el más delicado del capítulo, y en la pizarra se escribió dibujando \emph{dos caminos} sobre el mismo diagrama: uno en negro, que es la identidad de zigzag que se quiere probar, y otro en rojo, que es la identidad de zigzag de la hipótesis, con las flechas de naturalidad que los conectan. Trabajar en el modelo estricto, como se hace aquí, evita tener que arrastrar los asociadores en cada paso, pero oculta el trazado de los dos caminos; conviene tener el diagrama del pizarrón a la vista al seguir la sustitución.
\end{nota}
\end{proof}

\begin{ejer}[La otra identidad de zigzag]
Completar la demostración de la \myref{prop:simetria-dual} comprobando la identidad de zigzag restante.
\end{ejer}

\begin{sol}
Es el mismo cálculo con los papeles de $X$ y $X^\vee$ intercambiados. La composición que hay que verificar es
\[
X\xrightarrow{\ \lambda_X^{-1}\ }1\otimes X\xrightarrow{\ \mathrm{coev}'\otimes1_X\ }(X^\vee\otimes X)\otimes X\xrightarrow{\ \alpha_{X^\vee,X,X}\ }X^\vee\otimes(X\otimes X)\cdots
\]
y la sustitución de $\mathrm{coev}'=B_{X,X^\vee}\circ\mathrm{coev}$ y $\mathrm{ev}'=\mathrm{ev}\circ B_{X,X^\vee}$ reduce la expresión, en el modelo estricto, a
\[
(1_X\otimes\mathrm{ev}')\circ(\mathrm{coev}'\otimes1_X)=(1_X\otimes\mathrm{ev})\circ(1_X\otimes B_{X^\vee,X})\circ(B_{X,X^\vee}\otimes1_X)\circ(\mathrm{coev}\otimes1_X),
\]
que vuelve a ser la identidad de zigzag original con la trenza insertada dos veces; el hexágono las cancela y el resultado es $1_X$. Nótese que en esta dirección el orden en que aparecen las dos trenzas se invierte respecto de la anterior, que es justamente lo que hace que ambas se cancelen.
\end{sol}

%%%%%%%%%%%%%%%%%%%%%%%%%%%%%%%%%%%%%%%%%%%%%%%%%%%%%%%%%%%%%%%%%%%%%%%%%%%%%%%
\subsection{Unicidad del dual}
%%%%%%%%%%%%%%%%%%%%%%%%%%%%%%%%%%%%%%%%%%%%%%%%%%%%%%%%%%%%%%%%%%%%%%%%%%%%%%%

El dual de un objeto, cuando existe, es único en el sentido más débil posible: único salvo isomorfismo.

\begin{coro}\label{coro:unicidad-dual}
Si $X$ admite un dual, este es único salvo isomorfismo natural.
\end{coro}
\begin{proof}
Sean $X^\vee$ y $X^*$ dos duales de $X$. Por la \myref{prop:adjuncion}, los funtores $-\otimes X^\vee$ y $-\otimes X^*$ son ambos adjuntos derechos de $-\otimes X$. Como los adjuntos derechos son únicos salvo isomorfismo natural, existe un isomorfismo natural $-\otimes X^\vee\cong-\otimes X^*$. Evaluándolo en el objeto unidad se obtiene
\[
X^\vee\cong1\otimes X^\vee\cong1\otimes X^*\cong X^*,
\]
donde los isomorfismos de los extremos son las unidades monoidales.
\end{proof}

\begin{nota}
   La unicidad es «salvo isomorfismo» y no «a secas», y esa diferencia tendrá consecuencias en la sección siguiente: significa que \emph{no} hay una manera canónica de asignar a cada objeto su dual, sino solo una clase de isomorfismo. Es la misma situación que en álgebra homológica con los funtores derivados: el funtor $\mathrm{Tor}$ que se define al principio depende de las resoluciones proyectivas elegidas, y solo después se comprueba que resoluciones distintas dan funtores naturalmente isomorfos. En las categorías que aparecen en la práctica —y en particular en las de espectros— suele haber una elección canónica y el problema desaparece. (Este paralelismo con el álgebra homológica se explicó oralmente en clase: el profesor lo compara expresamente con «cuando dices que quieres hacer funtores derivados y tienes que escoger una resolución».)
\end{nota}

%%%%%%%%%%%%%%%%%%%%%%%%%%%%%%%%%%%%%%%%%%%%%%%%%%%%%%%%%%%%%%%%%%%%%%%%%%%%%%%
\subsection{El dual de un morfismo}
%%%%%%%%%%%%%%%%%%%%%%%%%%%%%%%%%%%%%%%%%%%%%%%%%%%%%%%%%%%%%%%%%%%%%%%%%%%%%%%

La \myref{coro:unicidad-dual} dice que tomar el dual no define un funtor: para definir un funtor habría que elegir un dual para cada objeto, y esa elección no es canónica. Lo que sí se puede hacer es fijar los duales y construir entonces el funtor.

La costumbre de trabajar con el dual de espacios vectoriales induce aquí una mala costumbre: la de pensar que el dual define siempre un funtor. En este contexto general no es así, pero casi. En el momento en que se fijan los duales, con los duales escogidos sí se tiene un funtor; la única libertad que queda es la elección, y esa elección es única salvo isomorfismo.

Supongamos fijados duales $X^\vee$ e $Y^\vee$ de $X$ e $Y$. La composición
\[
\hom_{\mathcal C}(X,Y)\cong\hom_{\mathcal C}(X\otimes Y^\vee,\ 1)\cong\hom_{\mathcal C}(Y^\vee\otimes X,\ 1)\cong\hom_{\mathcal C}(Y^\vee,\ X^\vee)
\]
es una biyección, donde el primer isomorfismo usa que $Y$ es un dual de $Y^\vee$ —por la \myref{prop:simetria-dual}— y el último que $X^\vee$ es un dual de $X$. La imagen de $f$ se denota $f^\vee$ y se llama el \emph{dual} de $f$.

Explícitamente, $f^\vee\colon Y^\vee\to X^\vee$ es la composición
\[
Y^\vee\xrightarrow{\ 1_{Y^\vee}\otimes\mathrm{coev}\ }Y^\vee\otimes X\otimes X^\vee\xrightarrow{\ 1_{Y^\vee}\otimes f\otimes1_{X^\vee}\ }Y^\vee\otimes Y\otimes X^\vee\xrightarrow{\ \mathrm{ev}\otimes1_{X^\vee}\ }1\otimes X^\vee\xrightarrow{\ \lambda_{X^\vee}\ }X^\vee .
\]

\begin{ejer}[El dual de la identidad y de una composición]
Comprobar que $1_X^\vee=1_{X^\vee}$ y que $(g\circ f)^\vee=f^\vee\circ g^\vee$, eligiendo en cada caso los duales de manera que las composiciones tengan sentido.
\end{ejer}

\begin{sol}
Para la identidad, la biyección $\hom(X,X)\to\hom(X^\vee,X^\vee)$ envía $1_X$ a $1_{X^\vee}$: la composición que define $1_X^\vee$ es
\[
X^\vee\xrightarrow{\ 1\otimes\mathrm{coev}\ }X^\vee\otimes X\otimes X^\vee\xrightarrow{\ 1\otimes1_X\otimes1\ }X^\vee\otimes X\otimes X^\vee\xrightarrow{\ \mathrm{ev}\otimes1\ }1\otimes X^\vee\xrightarrow{\ \lambda\ }X^\vee,
\]
que es la identidad de zigzag de la \myref{defi:dual} y por tanto vale $1_{X^\vee}$.

Para la composición, sean $f\colon X\to Y$ y $g\colon Y\to Z$ y fíjense duales $X^\vee$, $Y^\vee$, $Z^\vee$. La construcción de $(g\circ f)^\vee$ inserta $\mathrm{coev}$ de $X$ y evalúa con $\mathrm{ev}$ de $Z$; la de $f^\vee\circ g^\vee$ inserta $\mathrm{coev}$ de $Y$ en $g^\vee$ y evalúa con $\mathrm{ev}$ de $Y$ en $f^\vee$. La igualdad se sigue de la identidad de zigzag de $Y$, que dice que insertar $\mathrm{coev}$ y después $\mathrm{ev}$ en $Y$ es la identidad; el resto es reasociar, lo que el \myref{teo:mackane-kelly} permite hacer sin verificar nada.
\end{sol}

\begin{ejer}[El cuadrado que caracteriza a $f^\vee$]
Probar que el diagrama
\[
\begin{tikzcd}[column sep=large, row sep=large]
Y^\vee\otimes X \arrow[r,"1_{Y^\vee}\otimes f"] \arrow[d,"f^\vee\otimes1_X"'] & Y^\vee\otimes Y \arrow[d,"\mathrm{ev}_Y"]\\
X^\vee\otimes X \arrow[r,"\mathrm{ev}_X"'] & 1
\end{tikzcd}
\]
conmuta, y que la conmutatividad de este diagrama caracteriza a $f^\vee$.
\end{ejer}

\begin{sol}
La conmutatividad se comprueba sustituyendo la definición de $f^\vee$ y simplificando con las identidades de zigzag; es un cálculo diagramático del mismo tipo que los anteriores.

Para la caracterización, supongamos que $h\colon Y^\vee\to X^\vee$ hace conmutar el cuadrado. Entonces $\mathrm{ev}_X\circ(h\otimes1_X)=\mathrm{ev}_Y\circ(1_{Y^\vee}\otimes f)$, y componiendo ambos miembros por la derecha con $1_{Y^\vee}\otimes\mathrm{coev}_X$ y usando la identidad de zigzag de $X$ se obtiene $h=f^\vee$. Esto da una manera más conceptual de pensar el dual: $f^\vee$ es el único morfismo que hace conmutar el cuadrado de la evaluación, es decir, el único que respeta el «emparejamiento» entre un objeto y su dual.
\end{sol}

%%%%%%%%%%%%%%%%%%%%%%%%%%%%%%%%%%%%%%%%%%%%%%%%%%%%%%%%%%%%%%%%%%%%%%%%%%%%%%%
\section{Categorías monoidales cerradas y el hom interno}
%%%%%%%%%%%%%%%%%%%%%%%%%%%%%%%%%%%%%%%%%%%%%%%%%%%%%%%%%%%%%%%%%%%%%%%%%%%%%%%

%%%%%%%%%%%%%%%%%%%%%%%%%%%%%%%%%%%%%%%%%%%%%%%%%%%%%%%%%%%%%%%%%%%%%%%%%%%%%%%
\subsection{Categorías monoidales simétricas cerradas}
%%%%%%%%%%%%%%%%%%%%%%%%%%%%%%%%%%%%%%%%%%%%%%%%%%%%%%%%%%%%%%%%%%%%%%%%%%%%%%%

En una categoría monoidal simétrica general no todo objeto tiene dual, y el dual de un morfismo solo está definido tras elegir duales. Hay un contexto en el que las elecciones desaparecen porque existe un candidato canónico.

\begin{defi}\label{defi:cerrada}
Una categoría monoidal simétrica $\mathcal C$ es \emph{cerrada} si para todo objeto $X$ el funtor
\[
-\otimes X\colon\mathcal C\longrightarrow\mathcal C
\]
admite un adjunto derecho, que se denota $\ihom(X,-)$ y se llama el \emph{hom interno}.
\end{defi}

La adjunción se escribe, para todos $X,Y,Z$,
\[
\hom_{\mathcal C}(Y\otimes X,\ Z)\cong\hom_{\mathcal C}\bigl(Y,\ \ihom(X,Z)\bigr).
\]

\begin{ejem}\label{ejem:top-cerrada}
La categoría $\Top$ de espacios compactamente generados débilmente Hausdorff, con el producto, es monoidal simétrica cerrada. El hom interno es el espacio de mapeos con la topología compacto-abierta:
\[
\ihom(X,Z)=\Map(X,Z),
\]
y la adjunción es la conocida
\[
\hom_{\Top}(Y\times X,\ Z)\cong\hom_{\Top}\bigl(Y,\ \Map(X,Z)\bigr),
\]
que dice que una aplicación de dos variables es lo mismo que una aplicación de una variable al espacio de mapeos. La topología de $\Map(X,Z)$ es la que hace cierta esa biyección, y es la razón por la que se trabaja con espacios compactamente generados.
\end{ejem}

\begin{nota}
   El hom interno es un funtor en cada variable. Es covariante en la segunda por construcción, y contravariante en la primera; esto último es un ejercicio que se comprueba con la adjunción y la naturalidad. Al funtor de dos variables
   \[
   \ihom(-,-)\colon\mathcal C^{\mathrm{op}}\times\mathcal C\longrightarrow\mathcal C
   \]
   se le llama el \emph{hom interno de $\mathcal C$}. En las categorías de espectros que aparecerán más adelante el hom interno es una mejora del conjunto de morfismos, y su existencia es una de las cosas que se ganan al pasar a espectros.

   No hay una receta general para calcular el dual de un objeto, ni tampoco para construir su hom interno: cada categoría exige su propio método y no se dispone de una descripción que sirva para todas a la vez. En el caso de los espectros, la existencia de un hom interno se prueba con relativa facilidad, pero construirlo explícitamente es otra cosa.
\end{nota}

%%%%%%%%%%%%%%%%%%%%%%%%%%%%%%%%%%%%%%%%%%%%%%%%%%%%%%%%%%%%%%%%%%%%%%%%%%%%%%%
\subsection{El hom interno y el dual funcional}
%%%%%%%%%%%%%%%%%%%%%%%%%%%%%%%%%%%%%%%%%%%%%%%%%%%%%%%%%%%%%%%%%%%%%%%%%%%%%%%

La noción de hom interno está relacionada con la de dual por la unicidad de los adjuntos.

\begin{prop}\label{prop:hom-dual}
Si $X$ es dualizable con dual $X^\vee$, entonces existe un isomorfismo natural
\[
\ihom(X,-)\cong-\otimes X^\vee .
\]
En particular,
\[
X^\vee\cong1\otimes X^\vee\cong\ihom(X,1).
\]
\end{prop}
\begin{proof}
Por la \myref{prop:adjuncion}, $-\otimes X^\vee$ es adjunto derecho de $-\otimes X$; por la \myref{defi:cerrada}, $\ihom(X,-)$ también lo es. La unicidad de los adjuntos derechos da el isomorfismo natural. Evaluando en el objeto unidad y usando la unidad monoidal se obtiene la segunda afirmación.
\end{proof}

Lo interesante de la \myref{prop:hom-dual} es que la expresión $\ihom(X,1)$ tiene sentido aunque $X$ no sea dualizable, porque el hom interno existe para todo objeto en una categoría cerrada. Eso permite definir un candidato a dual en toda la categoría.

\begin{defi}\label{defi:dual-funcional}
Sea $\mathcal C$ una categoría monoidal simétrica cerrada. El funtor contravariante
\[
\ihom(-,1)\colon\mathcal C^{\mathrm{op}}\longrightarrow\mathcal C,
\]
y también cada uno de sus valores, se llama el \emph{dual funcional}. Se denota a veces
\[
DX\coloneqq\ihom(X,1).
\]
\end{defi}

\begin{nota}
   El dual funcional está definido sobre toda la categoría y es funtor sin necesidad de elegir nada: al ser el hom interno un funtor en cada variable, $\ihom(-,1)$ lo es. Sin embargo, \emph{no siempre es un dual}: la \myref{prop:hom-dual} dice que si $X$ es dualizable entonces $DX$ es un dual de $X$, pero el recíproco es falso. El ejemplo siguiente lo muestra.
\end{nota}

%%%%%%%%%%%%%%%%%%%%%%%%%%%%%%%%%%%%%%%%%%%%%%%%%%%%%%%%%%%%%%%%%%%%%%%%%%%%%%%
\subsection{Los objetos dualizables de los espacios vectoriales}
%%%%%%%%%%%%%%%%%%%%%%%%%%%%%%%%%%%%%%%%%%%%%%%%%%%%%%%%%%%%%%%%%%%%%%%%%%%%%%%

El ejemplo de los espacios vectoriales ilustra el fenómeno con toda claridad. Sea $k$ un cuerpo y consideremos la categoría $\Vect_k$ de \emph{todos} los espacios vectoriales sobre $k$, sin restricción de dimensión, con el producto tensorial $\otimes_k$.

\begin{prop}\label{prop:vect}
$(\Vect_k,\otimes_k,k)$ es una categoría monoidal simétrica cerrada, y sus objetos dualizables son exactamente los espacios vectoriales de dimensión finita.
\end{prop}
\begin{proof}
Que es simétrica monoidal es inmediato. Que es cerrada se sigue de la adjunción
\[
\Hom_k(A\otimes_k B,\ C)\cong\Hom_k\bigl(A,\ \Hom_k(B,C)\bigr),
\]
donde el hom interno de dos espacios vectoriales es el espacio de aplicaciones lineales, que tiene estructura de espacio vectorial sobre $k$.

Conviene precisar el estatus de los morfismos de la dualidad: dar una evaluación no quiere decir evaluar de verdad, sino que exista un morfismo formal de modo que unos diagramas conmuten. Lo mismo vale para la coevaluación.

Veamos que los dualizables son los de dimensión finita. Si $V$ es dualizable con dual $V^\vee$, la coevaluación es un morfismo
\[
\mathrm{coev}\colon k\longrightarrow V\otimes_k V^\vee,
\]
que está determinado por la imagen de $1$, y esa imagen es una suma finita
\[
1\longmapsto\sum_{i=1}^{m}v_i\otimes f_i
\]
con $v_i\in V$ y $f_i\in V^\vee$. Obsérvese que no se pide que los $v_i$ formen una base ni los $f_i$ una base dual: la suma finita es forzosa porque $1$ es un elemento de $V\otimes_k V^\vee$, y todo elemento de un producto tensorial es suma finita de tensores elementales.

En los cálculos que siguen se omiten los asociadores y las unidades, que no intervienen en el argumento.

Consideremos ahora la composición
\[
V^\vee\xrightarrow{\ 1\otimes\mathrm{coev}\ }V^\vee\otimes_k V\otimes_k V^\vee\xrightarrow{\ \mathrm{ev}\otimes1\ }k\otimes_k V^\vee\cong V^\vee .
\]
Sobre un elemento $f\in V^\vee$ actúa así:
\[
f\longmapsto\sum_{i=1}^{m}f\otimes v_i\otimes f_i\longmapsto\sum_{i=1}^{m}\mathrm{ev}(f\otimes v_i)\otimes f_i=\sum_{i=1}^{m}f(v_i)\,f_i,
\]
donde el último paso usa que el isomorfismo $k\otimes_k V^\vee\cong V^\vee$ está dado por multiplicar. Por la identidad de zigzag esta composición es la identidad de $V^\vee$, de modo que para todo $f\in V^\vee$ se tiene
\[
f=\sum_{i=1}^{m}f(v_i)\,f_i\in\langle f_1,\dots,f_m\rangle .
\]
Es decir, $V^\vee$ está contenido en el subespacio generado por los $m$ elementos fijos $f_1,\dots,f_m$, y por tanto es de dimensión finita; como $\dim V^\vee=\dim V$ cuando ambos son finitos, $V$ también lo es. Recíprocamente, si $V$ es de dimensión finita ya se sabe que $V^\vee=\Hom_k(V,k)$ es un dual de $V$ (véase el \myref{ejem:vect}).
\end{proof}

\begin{nota}
   Lo que el ejemplo ilustra es que el dual funcional $DV=\Hom_k(V,k)$ está definido para cualquier espacio vectorial, incluso de dimensión infinita, y que sin embargo solo es un dual en el sentido de categorías monoidales simétricas cuando $V$ es de dimensión finita. En dimensión infinita $DV$ es estrictamente más grande que $V$ y la coevaluación no puede existir, porque $1$ tendría que ir a una suma infinita.
\end{nota}

%%%%%%%%%%%%%%%%%%%%%%%%%%%%%%%%%%%%%%%%%%%%%%%%%%%%%%%%%%%%%%%%%%%%%%%%%%%%%%%
\subsection{El hom interno como mejora de los conjuntos de morfismos}
%%%%%%%%%%%%%%%%%%%%%%%%%%%%%%%%%%%%%%%%%%%%%%%%%%%%%%%%%%%%%%%%%%%%%%%%%%%%%%%

El hom interno no es solo un candidato a dual: es también la manera de dar estructura a los conjuntos de morfismos.

\begin{prop}\label{prop:hom-mor}
En una categoría monoidal simétrica cerrada existe una biyección natural
\[
\hom_{\mathcal C}(1,\ \ihom(X,Y))\cong\hom_{\mathcal C}(X,Y).
\]
\end{prop}
\begin{proof}
Se usa la adjunción que define el hom interno con $Y=1$:
\[
\hom_{\mathcal C}(1,\ \ihom(X,Y))\cong\hom_{\mathcal C}(1\otimes X,\ Y)\cong\hom_{\mathcal C}(X,Y),
\]
donde el último paso es la unidad monoidal.
\end{proof}

La \myref{prop:hom-mor} dice que el conjunto de morfismos $\hom_{\mathcal C}(X,Y)$ se recupera de $\ihom(X,Y)$ aplicando $\hom_{\mathcal C}(1,-)$, y que esa operación \emph{olvida} información. En este sentido el hom interno es una versión mejorada del conjunto de morfismos: contiene la misma información sobre los morfismos de $X$ a $Y$, pero además la estructura de objeto de $\mathcal C$.

\begin{ejem}\label{ejem:top-hom}
En $\Top$ el hom interno es $\Map(X,Y)$ \emph{como espacio}, mientras que $\hom_{\Top}(X,Y)$ es $\Map(X,Y)$ \emph{como conjunto}: al pasar de uno a otro se olvida la topología. La biyección de la \myref{prop:hom-mor} es aquí obvia, pues $\hom_{\Top}(*,\Map(X,Y))$ es el conjunto de puntos de $\Map(X,Y)$, es decir, el conjunto de aplicaciones de $X$ a $Y$.
\end{ejem}

\begin{ejem}\label{ejem:ch}
Sea $\Ch(\Ab)$ la categoría de complejos de cadenas de grupos abelianos, con el producto tensorial de complejos. Es monoidal simétrica cerrada, y el hom interno es el complejo de homomorfismos de complejos:
\[
\ihom(C_*,D_*)_n=\prod_{i\in\mathbb Z}\Hom_{\Ab}(C_i,\ D_{i+n}),
\]
con las diferenciales inducidas por las de $C_*$ y $D_*$, que llevan signos. En cambio, los morfismos de complejos son las familias $(f_i\colon C_i\to D_i)$ que conmutan con las diferenciales, es decir, tales que
\[
\partial'_i\circ f_i=f_{i-1}\circ\partial_i .
\]
La diferencia es visible: los morfismos de $\ihom(C_*,D_*)_n$ van de $C_i$ a $D_{i+n}$ y saltan de grado, mientras que los morfismos de complejos van de $C_i$ a $D_i$ y no saltan. La relación entre ambos es la esperada por la \myref{prop:hom-mor}: tomando $Z[0]$ como unidad monoidal,
\[
\hom_{\Ch(\Ab)}(Z[0],\ \ihom(C_*,D_*)_*)\cong H_0\bigl(\ihom(C_*,D_*)_*\bigr)\cong\hom_{\Ch(\Ab)}(C_*,D_*),
\]
que es el conocido cálculo del grupo de homología cero del complejo de homomorfismos.
\end{ejem}

\begin{nota}
   El \myref{ejem:ch} es el más ilustrativo de los tres: en los dos primeros el hom interno y el conjunto de morfismos están dados por el mismo conjunto subyacente, lo que oculta la diferencia, mientras que aquí el hom interno es un complejo con mucha más estructura. Conviene volver a este ejemplo cuando la notación del hom interno resulte confusa.
\end{nota}

%%%%%%%%%%%%%%%%%%%%%%%%%%%%%%%%%%%%%%%%%%%%%%%%%%%%%%%%%%%%%%%%%%%%%%%%%%%%%%%
\subsection{Propiedades del hom interno}
%%%%%%%%%%%%%%%%%%%%%%%%%%%%%%%%%%%%%%%%%%%%%%%%%%%%%%%%%%%%%%%%%%%%%%%%%%%%%%%

El hom interno satisface una lista de propiedades que son las análogas de las de los conjuntos de morfismos. Todas ellas se obtienen como adjuntos de morfismos evidentes.

\begin{prop}\label{prop:hom-propiedades}
En una categoría monoidal simétrica cerrada $\mathcal C$:
\begin{enumerate}[label = \textup{(\roman*)}]
\item \emph{Evaluación.} El morfismo
\[
X\otimes\ihom(X,Y)\xrightarrow{\ \mathrm{ev}\ }Y
\]
es el adjunto de la identidad de $\ihom(X,Y)$, es decir, corresponde a $1_{\ihom(X,Y)}$ bajo
\[
\hom_{\mathcal C}\bigl(X\otimes\ihom(X,Y),\ Y\bigr)\cong\hom_{\mathcal C}\bigl(\ihom(X,Y),\ \ihom(X,Y)\bigr).
\]
\item \emph{Composición.} Existe un morfismo
\[
\ihom(X,Y)\otimes\ihom(Y,Z)\longrightarrow\ihom(X,Z),
\]
adjunto de
\[
X\otimes\ihom(X,Y)\otimes\ihom(Y,Z)\xrightarrow{\ \mathrm{ev}\otimes1\ }Y\otimes\ihom(Y,Z)\xrightarrow{\ \mathrm{ev}\ }Z .
\]
\item \emph{Identidad.} Existe un morfismo
\[
1\xrightarrow{\ u\ }\ihom(X,X),
\]
adjunto de $X\otimes1\xrightarrow{\ \rho_X\ }X$. El morfismo $u$ es la \emph{unidad por la izquierda} para la composición del hom interno: el compuesto
\[
1\otimes\ihom(X,Y)\xrightarrow{\ u\otimes1\ }\ihom(X,X)\otimes\ihom(X,Y)\xrightarrow{\ \text{composición}\ }\ihom(X,Y)
\]
es la unidad $\lambda_{\ihom(X,Y)}$, como se comprueba adjuntando los dos miembros y usando que la composición del hom interno es el adjunto de $\mathrm{ev}\circ(\mathrm{ev}\otimes1)$. Este es el diagrama que se dibujó en la pizarra junto al apartado.
\item \emph{Ensamble.} Existe un morfismo
\[
\ihom(X,Y)\otimes Z\longrightarrow\ihom(X,\ Y\otimes Z),
\]
adjunto de
\[
X\otimes\ihom(X,Y)\otimes Z\xrightarrow{\ \mathrm{ev}\otimes1\ }Y\otimes Z .
\]
\item \emph{Tensor externo.} Existe un morfismo
\[
\ihom(X,Y)\otimes\ihom(W,Z)\longrightarrow\ihom(X\otimes W,\ Y\otimes Z),
\]
adjunto de
\[
X\otimes W\otimes\ihom(X,Y)\otimes\ihom(W,Z)\xrightarrow{\ \mathrm{ev}\otimes\mathrm{ev}\ }Y\otimes Z,
\]
donde el dominio se reordena para que las evaluaciones actúen sobre los pares correctos.
\end{enumerate}
\end{prop}

Todas estas propiedades son las análogas de las de los conjuntos de morfismos, y todas se prueban igual: se construye el morfismo del miembro derecho y se toma su adjunto. El \myref{teo:mackane-kelly} se encarga de los reordenamientos. Los dos últimos apartados admiten además una lectura común: el producto tensorial es un funtor monoidal, de modo que el comportamiento que tiene respecto de los morfismos lo tiene también respecto de los homs internos.

Falta en la lista la traducción de la asociatividad de los conjuntos de morfismos. Su análoga se obtiene con la misma composición del hom interno: componer tres factores usando la composición en los dos primeros y después en el tercero, o en el primero y después en los dos últimos, da el mismo resultado salvo el reordenamiento de los tres factores, y se comprueba del mismo modo que las anteriores.

\begin{prop}\label{prop:ensamble-iso}
Si $Z$ es dualizable, el morfismo de ensamble
\[
\ihom(X,Y)\otimes Z\longrightarrow\ihom(X,\ Y\otimes Z)
\]
es un isomorfismo.
\end{prop}
\begin{proof}[Idea de la demostración]
Se construye el inverso. Como $Z$ es dualizable, para dar un morfismo
\[
\ihom(X,\ Y\otimes Z)\longrightarrow\ihom(X,Y)\otimes Z
\]
basta, por adjunción, con dar
\[
\ihom(X,\ Y\otimes Z)\longrightarrow\ihom(Z^\vee,\ \ihom(X,Y)),
\]
y por la \myref{prop:hom-dual} el miembro derecho es naturalmente isomorfo a $Z^\vee\otimes\ihom(X,Y)$. Así, basta con exhibir un morfismo
\[
Z^\vee\otimes\ihom(X,\ Y\otimes Z)\longrightarrow\ihom(X,Y),
\]
que se obtiene componiendo la evaluación de $Z$ con la del hom interno. Es la hipótesis «$Z$ dualizable» la que habilita todas estas maniobras: cada vez que $Z$ pasa de un lado al otro de la adjunción se está usando su dual. Comprobar que este morfismo es el inverso requiere jugar con la naturalidad de la adjunción; no se hizo en clase.
\end{proof}

\begin{nota}
   El morfismo de ensamble aparece en muchas situaciones de álgebra y de topología, y a menudo interesa saber cuándo es un isomorfismo. La \myref{prop:ensamble-iso} da una condición suficiente: que $Z$ sea dualizable. El ejemplo clásico es el de los $R$-módulos, donde el ensamble
   \[
   \Hom_R(X,Y)\otimes_R Z\longrightarrow\Hom_R(X,\ Y\otimes_R Z)
   \]
   no es un isomorfismo en general pero sí lo es cuando $Z$ es un $R$-módulo libre finitamente generado, que es precisamente el caso dualizable. El mismo fenómeno reaparecerá con los espectros.
\end{nota}

El desarrollo completo de estas propiedades, con casi todas sus demostraciones, está en \cite{malkiewich}, la referencia principal del curso; allí algunas de ellas quedan como ejercicio.

%%%%%%%%%%%%%%%%%%%%%%%%%%%%%%%%%%%%%%%%%%%%%%%%%%%%%%%%%%%%%%%%%%%%%%%%%%%%%%%
\subsection{Categorías enriquecidas sobre sí mismas}
%%%%%%%%%%%%%%%%%%%%%%%%%%%%%%%%%%%%%%%%%%%%%%%%%%%%%%%%%%%%%%%%%%%%%%%%%%%%%%%

La \myref{prop:hom-mor} tiene una lectura estructural. Una categoría enriquecida sobre una categoría monoidal $\mathcal V$ es una categoría en la que los morfismos entre dos objetos son objetos de $\mathcal V$ y la composición es un morfismo de $\mathcal V$. En una categoría monoidal simétrica cerrada $\mathcal C$, los homs internos son objetos de $\mathcal C$ y las propiedades del hom interno dan precisamente eso. En todo el capítulo se trabaja con categorías localmente pequeñas, de modo que $\hom_{\mathcal C}(X,Y)$ es un conjunto.

\begin{prop}\label{prop:enriquecida}
Sea $\mathcal C$ una categoría monoidal simétrica cerrada. Entonces $\mathcal C$ es una categoría enriquecida sobre sí misma, con
\[
\mathcal C(X,Y)\coloneqq\ihom(X,Y)
\]
y con la composición dada por el morfismo del hom interno construido en la \myref{prop:hom-propiedades}.
\end{prop}

La categoría original se recupera de la enriquecida aplicando $\hom_{\mathcal C}(1,-)$ a los homs internos, por la \myref{prop:hom-mor}, y esa recuperación es compatible con la composición. Hay que tener cuidado en no confundir las dos categorías: la enriquecida es una categoría distinta, con más estructura, y la original es su «imagen débil» bajo $\hom_{\mathcal C}(1,-)$.

\begin{nota}
   El ejemplo de los complejos de cadenas lo aclara: la categoría de complejos enriquecida sobre sí misma tiene como objetos de morfismos los complejos $\ihom(C_*,D_*)_*$, y al aplicar $H_0$ se recuperan los morfismos de complejos. La versión enriquecida conserva además toda la homología superior, que la versión ordinaria no ve.

   El mismo cuidado hay que tener con la categoría de espectros: al enriquecerla sobre sí misma tomando como objetos de morfismos los espectros $\ihom(X,Y)$, la categoría que resulta se sigue llamando categoría de espectros, pero no es la categoría de espectros de partida.

   Una categoría enriquecida sobre $\mathcal C$ puede a su vez ser monoidal simétrica, pero la simetría hay que pedirla en el sentido enriquecido, lo que requiere que $\mathcal C$ sea monoidal simétrica; aquí lo es, ya que nos estamos enriqueciendo sobre $\mathcal C$ misma.
\end{nota}

%%%%%%%%%%%%%%%%%%%%%%%%%%%%%%%%%%%%%%%%%%%%%%%%%%%%%%%%%%%%%%%%%%%%%%%%%%%%%%%
\section{La dualidad de Spanier--Whitehead sobre las esferas}
%%%%%%%%%%%%%%%%%%%%%%%%%%%%%%%%%%%%%%%%%%%%%%%%%%%%%%%%%%%%%%%%%%%%%%%%%%%%%%%

%%%%%%%%%%%%%%%%%%%%%%%%%%%%%%%%%%%%%%%%%%%%%%%%%%%%%%%%%%%%%%%%%%%%%%%%%%%%%%%
\subsection{La dualidad sobre las esferas}
%%%%%%%%%%%%%%%%%%%%%%%%%%%%%%%%%%%%%%%%%%%%%%%%%%%%%%%%%%%%%%%%%%%%%%%%%%%%%%%

Se vuelve ahora al contexto estable. Recordemos que $(\SW,\wedge)$ es una categoría monoidal simétrica, con unidad monoidal $(S^0,0)$. Lo que \emph{no} se sabe todavía es si es cerrada: no hay por ahora un adjunto derecho conocido para $-\wedge X$, y por tanto tampoco un candidato canónico a dual. Pero las consecuencias generales ya dan bastante.

En el contexto estable se usa la notación $X^*$ para el dual, y a veces $DX$. El hom interno, cuando exista, tiene que ser un objeto de la misma categoría, es decir, un espectro, y se le suele denotar $F(X,Y)$. Empecemos por la unidad.

\begin{prop}\label{prop:s0-dual}
La unidad monoidal $(S^0,0)$ es dualizable y es su propio dual: $(S^0,0)^*=(S^0,0)$.
\end{prop}
\begin{proof}
Es el \myref{coro:unidad-autodual} aplicado a $\SW$.
\end{proof}

A continuación, las esferas. Recordemos la propiedad de desplazamiento de índices, que es la forma en que se expresa la estabilidad:
\[
\set*{(A,m+r),\ (B,n)}\cong\set*{(A,m),\ (B,n-r)}
\]
naturalmente en los objetos $(A,m)$ y $(B,n)$. Como se hizo en la pizarra, se mantienen ambos objetos arbitrarios y se reescriben los desplazamientos como productos smash:
\[
\begin{aligned}
\set*{(A,m)\wedge(S^0,r),\ (B,n)}
&\cong\set*{(A,m+r),\ (B,n)}\\
&\cong\set*{(A,m),\ (B,n-r)}\\
&\cong\set*{(A,m),\ (B,n)\wedge(S^0,-r)}.
\end{aligned}
\]
Estas biyecciones son naturales en las dos variables, por lo que expresan la adjunción
\[
-\wedge(S^0,r)\ \dashv\ -\wedge(S^0,-r).
\]
Por la \myref{prop:adjuncion}, $(S^0,-r)$ es un dual de $(S^0,r)$. Escribiendo $S^r=(S^0,r)$ para todo $r\in\mathbb Z$, se obtiene
\[
S^{-r}=(S^r)^* .
\]

\begin{ejer}[Evaluación y coevaluación de las esferas]
Encontrar explícitamente los morfismos de evaluación y coevaluación que exhiben $S^{-r}$ como dual de $S^r$.
\end{ejer}

\begin{sol}
Se trabaja con los objetos $S^r=(S^0,r)$, $S^{-r}=(S^0,-r)$ y $S=(S^0,0)$ de $\SW$. Por la definición del producto smash,
\[
S^{-r}\wedge S^r=(S^0\wedge S^0,0),
\qquad
S^r\wedge S^{-r}=(S^0\wedge S^0,0).
\]
Sea $u\colon S^0\wedge S^0\to S^0$ el homeomorfismo punteado que envía el único punto no base al único punto no base. Se eligen
\[
\mathrm{ev}=[u]\colon S^{-r}\wedge S^r\longrightarrow S,
\qquad
\mathrm{coev}=[u^{-1}]\colon S\longrightarrow S^r\wedge S^{-r}.
\]
Son morfismos explícitos: en cualquier nivel $N\geq0$ están representados por $\Sigma^N u$ y $\Sigma^N(u^{-1})$, respectivamente. Identificando $\Sigma^N(S^0\wedge S^0)$ con $S^N$, ambos representantes son $\id_{S^N}$. Esto vale para cualquier entero $r$, pues los desplazamientos $-r+r$ y $r-r$ se anulan antes de escoger representantes.

Queda verificar los zigzags. Abreviemos $E=S^r$ y $D=S^{-r}$. Las dos composiciones que exige la \myref{defi:dual} son
\[
\begin{aligned}
&\rho_E\circ(1_E\wedge\mathrm{ev})\circ\alpha_{E,D,E}
\circ(\mathrm{coev}\wedge1_E)\circ\lambda_E^{-1},\\
&\lambda_D\circ(\mathrm{ev}\wedge1_D)\circ\alpha^{-1}_{D,E,D}
\circ(1_D\wedge\mathrm{coev})\circ\rho_D^{-1}.
\end{aligned}
\]
Si $e$ denota el punto no base de $S^0$, el primer compuesto inserta dos factores $e$ a la izquierda y elimina los dos de la derecha; el segundo los inserta a la derecha y elimina los dos de la izquierda. En ambos casos el punto base permanece fijo y el punto no base recorre
\[
e\longmapsto[e,e,e]\longmapsto e.
\]
Los asociadores solo cambian los paréntesis y los unitores eliminan un factor $S^0$. No se intercambian factores. Además, $u$ y $u^{-1}$ tienen desplazamiento cero tanto en la fuente como en el blanco, de modo que los signos de la definición del smash se cancelan al tensarlos con una identidad. Por tanto, en un nivel suficientemente alto los dos compuestos son las identidades de $S^{N+r}$ y $S^{N-r}$, respectivamente. Sus clases son $1_E$ y $1_D$, como se quería.
\end{sol}

\begin{nota}
   El ejercicio se dejó planteado en clase; la solución anterior es un desarrollo añadido. La evaluación elegida es la identificación estable de grado uno, no un colapso de un espacio de dimensión negativa.
\end{nota}

El resultado siguiente permite trasladar la dualidad de un objeto a sus suspensiones.

\begin{lema}\label{lema:suspension-dual}
Si $(Y,m)$ es un dual de $(X,n)$, entonces $(Y,m-k)$ es un dual de $(X,n+k)$ para todo $k\in\mathbb Z$.
\end{lema}
\begin{proof}
Por el \myref{coro:dual-tensor} y la \myref{prop:s0-dual},
\[
(X,n+k)=(X,n)\wedge(S^0,k),
\]
y como $(S^0,k)^*=(S^0,-k)$ por el párrafo anterior, el dual de un producto es el producto de los duales en el orden invertido:
\[
(X,n+k)^*=(S^0,k)^*\wedge(X,n)^*=(S^0,-k)\wedge(Y,m)=(Y,m-k). \qedhere
\]
\end{proof}

%%%%%%%%%%%%%%%%%%%%%%%%%%%%%%%%%%%%%%%%%%%%%%%%%%%%%%%%%%%%%%%%%%%%%%%%%%%%%%%
\subsection{Un criterio para la dualidad en \texorpdfstring{$\SW$}{SW}}
%%%%%%%%%%%%%%%%%%%%%%%%%%%%%%%%%%%%%%%%%%%%%%%%%%%%%%%%%%%%%%%%%%%%%%%%%%%%%%%

El objetivo de esta parte es demostrar que \emph{todo} objeto de $\SW$ es dualizable. Por el \myref{lema:suspension-dual} es suficiente probarlo para los objetos de la forma $(X,0)$, y el resultado siguiente reduce la comprobación a una condición sobre las esferas.

\begin{lema}\label{lema:criterio-dual}
Sea $X$ un objeto de $\SW$ y sea $X^*$ otro objeto. Entonces $X^*$ es un dual de $X$ si y solo si existe un morfismo
\[
\mathrm{ev}\colon X^*\wedge X\longrightarrow S,
\qquad
S\coloneqq(S^0,0),
\]
tal que las dos aplicaciones
\[
\set*{\Sigma^r S,\ X^*}\longrightarrow\set*{\Sigma^r X,\ S},
\qquad
\set*{\Sigma^r S,\ X}\longrightarrow\set*{\Sigma^r X^*,\ S}
\]
definidas por
\[
f\longmapsto\mathrm{ev}\circ(f\wedge1_X),
\qquad
f\longmapsto\mathrm{ev}\circ B_{X,X^*}\circ(f\wedge1_{X^*})
\]
son biyecciones para todo $r\in\mathbb Z$.
\end{lema}

\begin{proof}
\emph{Si $X^*$ es el dual de $X$, las dos aplicaciones son biyecciones.} Es la adjunción de la \myref{prop:adjuncion} escrita con $A=\Sigma^rS$ y $B=S$:
\[
\set*{\Sigma^rS,\ B\wedge X^*}\xrightarrow{\ \cong\ }\set*{\Sigma^rS\wedge X,\ B},
\qquad
f\longmapsto\mathrm{ev}\circ(f\wedge1_X),
\]
donde $B$ es cualquier objeto. Tomando $B=S$ se obtiene la primera biyección. La segunda es la misma adjunción con $X$ y $X^*$ intercambiados, usando la \myref{prop:simetria-dual} y que la evaluación correspondiente es $\mathrm{ev}'=\mathrm{ev}\circ B_{X,X^*}$, con la trenza $B_{X,X^*}\colon X\wedge X^*\to X^*\wedge X$.

\emph{Recíprocamente.} Supongamos que las biyecciones se cumplen y sea $A$ un CW-complejo finito. Se prueba que
\[
\set*{A,\ X^*}\xrightarrow{\ \cong\ }\set*{A\wedge X,\ S}
\]
por inducción sobre el número de celdas de $A$. Sea
\[
D\xrightarrow{\ f\ }B\xrightarrow{\ g\ }C\xrightarrow{\ h\ }\Sigma D
\]
una sucesión cofibrada de CW-complejos, donde $C$ se obtiene de $B$ pegando una celda a lo largo de $D=S^m$. Aplicando los dos funtores ---ambos contravariantes en la primera variable--- se obtiene el diagrama
\[
\begin{tikzcd}[column sep=small, row sep=large]
\set*{\Sigma B,\ X^*} \arrow[r] \arrow[d,"\cong"'] & \set*{\Sigma D,\ X^*} \arrow[r] \arrow[d,"\cong"'] & \set*{C,\ X^*} \arrow[r] \arrow[d,dashed] & \set*{B,\ X^*} \arrow[r] \arrow[d,"\cong"'] & \set*{D,\ X^*} \arrow[d,"\cong"']\\
\set*{\Sigma B\wedge X,\ S} \arrow[r] & \set*{\Sigma D\wedge X,\ S} \arrow[r] & \set*{C\wedge X,\ S} \arrow[r] & \set*{B\wedge X,\ S} \arrow[r] & \set*{D\wedge X,\ S}
\end{tikzcd}
\]
Las filas son exactas por la \myref{prop:exacta}, y están escritas en el orden de los morfismos inducidos: como los dos funtores son contravariantes, ese orden es el \emph{inverso} del de la sucesión cofibrada $D\to B\to C\to\Sigma D\to\Sigma B$. Las flechas verticales de los extremos son isomorfismos: las de $D$ y de $\Sigma D$ porque $D$ y $\Sigma D$ son esferas, y en ese caso la biyección es la hipótesis; las de $B$ y $\Sigma B$ por hipótesis de inducción, ya que $B$ tiene menos celdas que $C$ y $\Sigma B$ tiene las mismas que $B$. Por el lema de los cinco, la flecha central también es un isomorfismo. Como los CW-complejos finitos se construyen con un número finito de pasos de este tipo, la biyección vale para todo CW-complejo finito $A$.

\emph{Paso a todos los objetos.} Sea ahora $A$ un objeto arbitrario de $\SW$ y $n$ tal que $\Sigma^nA$ sea un CW-complejo finito. Aplicando lo anterior a $\Sigma^nA$ y a los objetos $\Sigma^nX^*$ y $\Sigma^{-n}X$,
\[
\set*{\Sigma^nA,\ \Sigma^nX^*}\cong\set*{\Sigma^nA\wedge\Sigma^{-n}X,\ S},
\]
donde la evaluación que hace de $\Sigma^nX^*$ un dual de $\Sigma^{-n}X$ se construye explícitamente reordenando los factores y aplicando la evaluación original:
\[
\Sigma^nX^*\wedge\Sigma^{-n}X\xrightarrow{\ \cong\ }\mathbb S^n\wedge X^*\wedge\mathbb S^{-n}\wedge X\longrightarrow\mathbb S^n\wedge\mathbb S^{-n}\wedge X^*\wedge X\xrightarrow{\ \mathrm{ev}\wedge\mathrm{ev}\ }\mathbb S,
\]
que cumple las mismas propiedades que la de partida. Por la propiedad de desplazamiento ambos miembros se reescriben como
\[
\set*{A,\ X^*}\cong\set*{A\wedge X,\ S},
\]
con lo que la biyección vale para todo $A$. Todas las biyecciones que aparecen en el argumento son naturales, de modo que la biyección final es natural en $A$.

\emph{Segunda entrada.} Repitiendo el mismo argumento en la segunda variable se obtiene una biyección
\[
\set*{A,\ B\wedge X^*}\cong\set*{A\wedge X,\ B}
\]
natural en $A$ y en $B$.

\emph{Conclusión.} La biyección anterior dice que $-\wedge X$ es adjunto izquierdo de $-\wedge X^*$. Por la \myref{prop:adjuncion}, esto es exactamente decir que $X^*$ es un dual de $X$.
\end{proof}

\begin{nota}
   El lema reduce un problema sobre toda la categoría a un problema sobre las esferas, y la razón es estructural: todo objeto de $\SW$ se obtiene de $S^0$ mediante una cantidad finita de suspensiones y de sucesiones cofibradas, y las sucesiones cofibradas dan lugar a sucesiones exactas en cualquiera de las dos entradas. Por eso basta con tener la adjunción para las esferas y propagarla con el lema de los cinco.

   Lo que falta para cerrar el resultado general es el mismo argumento pero partiendo de una sucesión cofibrada con los tres términos arbitrarios: si dos de ellos son dualizables, el tercero también. Eso es más delicado que lo anterior, porque no hay una evaluación natural predefinida para el tercer objeto y hay que construirla, con detalles de topología de conjuntos. Se hará en la clase siguiente.
\end{nota}

\begin{nota}
   Los objetos dualizables no determinan por sí solos toda la categoría. Con los dualizables de la categoría homotópica estable no se obtienen todas las teorías de cohomología, pero sí todo lo que se detecta mediante condiciones de finitud, es decir, lo que concierne a los objetos compactos. La categoría no se recupera a partir de ellos sin más, porque siempre se pueden añadir objetos superfluos que no intervienen en esa parte.

   El paralelismo con los espacios vectoriales es instructivo: los espacios de dimensión infinita son colímites de los de dimensión finita y en ese sentido se recuperan de ellos, pero hay cosas que se pueden hacer con un espacio de dimensión infinita que no se reducen a lo que ocurre en dimensión finita. En la categoría derivada ocurre lo mismo con los complejos perfectos: hay propiedades de la categoría derivada que no provienen de ellos, y no hay una equivalencia entre la categoría derivada y la subcategoría plena de complejos perfectos. Lo que sí cabe esperar es que la categoría se genere a partir de los objetos dualizables mediante ciertas operaciones.
\end{nota}
